%------------------------------------------------------------------------------%
\begin{question}
  Determine a equação geral do plano que passa pelos pontos
  \[
    A = (1, 0, 2)
  \]
  \[
    B = (3, 1, 1)
  \]
  \[
    C = (0, 2, 3)
  \]
\end{question}

%------------------------------------------------------------------------------%
\begin{resolution}{Solução}
  Construímos dois vetores contidos no plano

  \pause
  \[
    \vec{v}
    \pause
    = \overrightarrow{AB}
    \pause
    = B - A
    \pause
    = (3, 1, 1) - (1, 0, 2)
    \pause
    = (2, 1, -1)
  \]

  \pause
  \[
    \vec{w}
    \pause
    = \overrightarrow{AC}
    \pause
    = C - A
    \pause
    = (0, 2, 3) - (1, 0, 2)
    \pause
    = (-1, 2, 1)
  \]
\end{resolution}

%------------------------------------------------------------------------------%
\begin{resolution}{Solução}
  \onslide<+->{Um vetor normal é}
  \begin{align*}
    \onslide<+->{\vec{n}}
    \onslide<+->{& = \vec{v} \times \vec{w} \\}
    \onslide<+->{& = \begin{vmatrix}
          \vec{i} & \vec{j} & \vec{k} \\
          2       & 1       & -1      \\
          -1      & 2       & 1
        \end{vmatrix}}
    \\
    \onslide<+->{& = (3, -1 ,5)}
  \end{align*}
\end{resolution}

%------------------------------------------------------------------------------%
\begin{resolution}{Solução}
  \onslide<+->{Usando o ponto $A=(1,0,2)$}
  \begin{align*}
    \onslide<+->{\vec{n} \cdot \overrightarrow{AP} & = 0 \\}
    \onslide<+->{(3, -1 ,5) \cdot (x-1, y-0, z-2)  & = 0 \\}
    \onslide<+->{3(x-1)-y+5(z-2)                   & = 0 \\}
    \onslide<+->{3x-3-y+5z-10                      & = 0 \\}
    \onslide<+->{3x-y+5z-13                        & = 0}
  \end{align*}
\end{resolution}

%------------------------------------------------------------------------------%
\endinput
